\documentclass[a4paper,12pt]{article}
\usepackage[utf8]{inputenc}
\usepackage[T1]{fontenc}
\usepackage[hungarian]{babel}
\usepackage{graphicx}
\usepackage{geometry}
\geometry{a4paper,
		     tmargin = 35mm, 
		     lmargin = 25mm,
		     rmargin = 30mm,
		     bmargin = 30mm}
\usepackage{mathtools}
\usepackage{amsmath}
\usepackage{color}
\usepackage{setspace}
\usepackage{amsmath,amssymb}
\usepackage{float}
\usepackage{hyperref}

\usepackage{indentfirst}
\usepackage{subfig}

\usepackage{siunitx}

\renewcommand\thesection{\Roman{section}}

\begin{document}

\linespread{1.25}

\begin{titlepage}

	\centering
	\includegraphics[width=0.66\textwidth]{elte.jpg}\par\vspace{1cm}
	{\scshape\LARGE ELTE TTK \par}
	\vspace{3cm}
	{\scshape\Large Röntgendiffrakció \par}
	\vspace{1cm}
	{\large\itshape Olar Alex\par}
	\vspace{3cm}
	{\large 2018 \par}
	
\end{titlepage}

\tableofcontents

\newpage

\section{Elméleti összefoglaló, mérési eszközök}

\vspace{5mm}

\par A mérés során 4 feladatot végeztünk el. Ennek során egy vas-alumínium ötvözetet vizsgáltunk először, melynek a rácsszerkezetét kellett meghatároznunk, majd egy alumínium-magnézium ötvözetre tértünk át, amiben a magnézium arányát kellett meghatároznunk megváltozott rácstávolság alapján. Majd egy ismeretlen minta (szén-pasztilla) csúcskiszélesedése alapján meg lehetett határozni az ismeretlen fázist. Végül egy otthoni feladatot kaptunk a rácsparaméter otthoni meghatározására.

\vspace{5mm}

\section{Kiértékelés}

\vspace{5mm}

\subsection{Rácstípus maghatározása $Fe-Al$ ötvözet}

\vspace{5mm}

\par A reflexiók indexelésével meghatároztuk a difraktorgramból, hogy az ötvözet $BCC$ rács. Ezt a következőképpen kiviteleztük:

\begin{itemize}
\item meghatároztuk a Bragg-egyenlet alapján, ($2d_{hkl}sin\vartheta = \lambda$), a $d_{hkl}$ értékeket, majd leosztva az első diffrakciós csúcshoz tartozó $d_{hkl}$ értékkel az egész sort, kaptuk a közelítő $N$ értékeket
\end{itemize} 

\vspace{5mm}

\begin{center}
\begin{tabular}{|c|c|c|}
\hline
$2\vartheta$ & $d_{hkl}$ [$\si{\angstrom}$] & N \\
\hline
30.7 &2.9121 & 1 \\
\hline
44.0 & 2.058 & 2.003 \\
\hline
54.6 & 1.681 & 3.002 \\
\hline
64.0 & 1.455 & 4.007 \\
\hline
72.6 & 1.302 & 5.002 \\
\hline
80.9 & 1.188 & 6.007 \\
\hline
97.1 & 1.029 & 8.017 \\
\hline
105.3 & 0.97 & 9.018 \\
\hline
113.8 & 0.92 & 10.015 \\
\hline
123.0 & 0.877 & 11.022 \\
\hline
133.2 & 0.84 & 12.02 \\
\hline
\end{tabular}
\end{center}

\vspace{5mm}

\par Ebből jól látható, hogy jó közelítéssel egész $N$ értékeket kaptunk, amiből látható, hogy ez egy $BCC$ rács, hiszen:

\begin{equation*}
\frac{1}{d_{hkl}^{2}} = \frac{1}{a^{2}}N = \frac{1}{a^{2}}(h^{2} + k^{2} + l^{2})
\end{equation*}

\par összefüggést csak köbös rácsra lehet alkalmazni.

\subsection{Ötvözőkoncentráció meghatározása}

\vspace{5mm}

\par Az alumínium magnézium szennyezettségét kellett meghatározni diffrakcióból. A kapott diffrakciós görbén  

\end{document}